\chapter{Euclidean Spaces}\label{ch:b1:euclid}

Adding an inner product to a real vector space buys the geometric
notions --- lengths, angles, orthogonality, distances --- and one
theorem that towers over the chapter: every subspace admits an
orthogonal projection, computable by Gram--Schmidt, realizing the
shortest distance. The plane isometries close the chapter and the
year's geometry.

Throughout, $E$ is a \emph{real} vector space.

\section{Inner products}

\begin{definition}\label{def:b1:euclid:def}
An \emph{inner product}\index{inner product} on $E$ is a map
$\langle\cdot,\cdot\rangle \colon E \times E \to \R$ that is
bilinear, symmetric, and positive definite ($\langle x, x\rangle >
0$ for $x \neq 0$). A finite-dimensional space so equipped is a
\emph{Euclidean space}\index{Euclidean space}. The \emph{norm} of
$x$ is $\norm{x} = \sqrt{\langle x, x\rangle}$, and $d(x, y) =
\norm{x - y}$.
\end{definition}

\begin{example}\label{ex:b1:euclid:examples}
On $\R^n$: the canonical product $\langle x, y\rangle = \sum x_i
y_i$. On $C(\intcc{a}{b})$: $\langle f, g \rangle = \int_a^b fg$
(positive definiteness is \cref{thm:b1:integration:props}~(4)). On
$\R_n[X]$: $\langle P, Q\rangle = \int_0^1 PQ$, or $\sum_{i}
P(x_i)Q(x_i)$ over $n+1$ distinct points.
\end{example}

\begin{theorem}[Cauchy--Schwarz; norm properties]\label{thm:b1:euclid:cs}
For all $x, y \in E$:\index{Cauchy--Schwarz inequality}
\[
\abs{\langle x, y\rangle} \leq \norm x\, \norm y ,
\]
with equality iff $x, y$ are proportional. Consequently $\norm\cdot$
satisfies the triangle inequality $\norm{x + y} \leq \norm x + \norm
y$ (and $\norm{\lambda x} = \abs\lambda \norm x$, $\norm x = 0 \iff
x = 0$). Moreover:
\[
\norm{x+y}^2 = \norm x^2 + 2\langle x, y\rangle + \norm y^2,
\qquad
\langle x, y \rangle = \tfrac14\bigl(\norm{x+y}^2 -
\norm{x-y}^2\bigr) .
\]
\end{theorem}

\begin{proof}
If $y = 0$, everything is trivial. Otherwise the quadratic $t
\mapsto \norm{x + ty}^2 = \norm x^2 + 2t\langle x, y\rangle + t^2
\norm y^2$ is $\geq 0$ for all $t$: its discriminant is $\leq 0$,
which is Cauchy--Schwarz; equality means a double root $t_0$, i.e.\
$x + t_0 y = 0$ (definiteness): proportionality. Triangle
inequality: expand $\norm{x+y}^2$ and bound the cross term by
Cauchy--Schwarz to get $(\norm x + \norm y)^2$. The last two
identities are direct expansions (the second, the
\emph{polarization identity}, recovers the product from the norm).
\end{proof}

\section{Orthogonality}

\begin{definition}\label{def:b1:euclid:orthogonal}
$x \perp y$ when $\langle x, y \rangle = 0$. A family is
\emph{orthogonal} when its vectors are pairwise orthogonal,
\emph{orthonormal}\index{orthonormal family} when moreover each has
norm $1$. The \emph{orthogonal complement} of a subspace $F$ is
\[
F^{\perp} = \{x \in E : \forall y \in F,\ \langle x, y\rangle =
0\},
\]
a subspace of $E$.\index{orthogonal complement}
\end{definition}

\begin{proposition}\label{prop:b1:euclid:orthfree}
(Pythagoras\index{Pythagorean theorem}) If $x \perp y$ then
$\norm{x+y}^2 = \norm x^2 + \norm y^2$. An orthogonal family of
nonzero vectors is free. In an orthonormal basis $(e_1, \dots,
e_n)$, coordinates and products are
\[
x = \sum_{i} \langle x, e_i\rangle\, e_i,
\qquad
\langle x, y \rangle = \sum_i \langle x, e_i\rangle \langle y,
e_i\rangle,
\qquad
\norm x^2 = \sum_i \langle x, e_i\rangle^2 .
\]
\end{proposition}

\begin{proof}
Pythagoras: expand. Freeness: take $\langle\,\cdot\,, x_j\rangle$
of a null combination: $\lambda_j \norm{x_j}^2 = 0$. Coordinates:
write $x = \sum \lambda_i e_i$ and take the product with $e_j$:
$\lambda_j = \langle x, e_j\rangle$; the two formulas follow by
bilinearity.
\end{proof}

\begin{theorem}[Gram--Schmidt]\label{thm:b1:euclid:gramschmidt}
Every Euclidean space has orthonormal bases. Explicitly, from any
basis $(v_1, \dots, v_n)$, the recipe\index{Gram--Schmidt process}
\[
w_k = v_k - \sum_{i=1}^{k-1} \langle v_k, e_i\rangle\, e_i ,
\qquad
e_k = \frac{w_k}{\norm{w_k}}
\]
produces an orthonormal basis $(e_1, \dots, e_n)$ with
$\operatorname{Vect}(e_1, \dots, e_k) = \operatorname{Vect}(v_1,
\dots, v_k)$ for every $k$.
\end{theorem}

\begin{proof}
Induction on $k$. Assuming $(e_1, \dots, e_{k-1})$ orthonormal
spanning $\operatorname{Vect}(v_1, \dots, v_{k-1})$: the vector
$w_k$ is orthogonal to each $e_j$ ($j < k$) by construction
($\langle w_k, e_j \rangle = \langle v_k, e_j\rangle - \langle v_k,
e_j\rangle$), and $w_k \neq 0$ since $v_k \notin
\operatorname{Vect}(v_1, \dots, v_{k-1})$. Normalizing keeps
orthogonality; the span statement holds since $e_k$ is a
combination of $v_k$ and earlier $e_i$'s, invertibly.
\end{proof}

\begin{theorem}[Orthogonal projection]\label{thm:b1:euclid:projection}
Let $F$ be a subspace of the Euclidean space $E$. Then
\[
E = F \oplus F^{\perp},
\]
and the associated projection $p_F$ onto $F$ (the \emph{orthogonal
projection}\index{orthogonal projection}) is given, in any
orthonormal basis $(e_1, \dots, e_k)$ of $F$, by $p_F(x) = \sum_i
\langle x, e_i\rangle e_i$. It realizes the distance to $F$:
for all $y \in F$,
\[
\norm{x - p_F(x)} \leq \norm{x - y},
\]
with equality only for $y = p_F(x)$; one writes $d(x, F) = \norm{x -
p_F(x)}$.
\end{theorem}

\begin{proof}
Take an orthonormal basis $(e_i)_{i \leq k}$ of $F$
(\cref{thm:b1:euclid:gramschmidt} inside $F$) and set $\pi(x) =
\sum \langle x, e_i\rangle e_i \in F$. Then $x - \pi(x) \perp e_j$
for each $j$ (same cancellation as above), hence $x - \pi(x) \in
F^\perp$: $E = F + F^\perp$. And $F \cap F^\perp = \{0\}$: such a
vector satisfies $\langle x, x\rangle = 0$. So the sum is direct and
$\pi = p_F$.

Distance: for $y \in F$, decompose $x - y = (x - p_F(x)) + (p_F(x) -
y)$, orthogonal pieces ($F^\perp$ and $F$); Pythagoras:
\[
\norm{x - y}^2 = \norm{x - p_F(x)}^2 + \norm{p_F(x) - y}^2
\geq \norm{x - p_F(x)}^2,
\]
equality iff $y = p_F(x)$.
\end{proof}

\begin{example}[Best quadratic approximation]\label{ex:b1:euclid:bestapprox}
In $C(\intcc{0}{1})$ with $\langle f, g\rangle = \int_0^1 fg$, the
polynomial of degree $\leq 1$ closest to $f(x) = x^2$ in the
associated (quadratic mean) distance is $p_F(f)$ where $F = \R_1[X]$.
Gram--Schmidt on $(1, X)$: $e_1 = 1$, $w_2 = X - \frac12$,
$\norm{w_2}^2 = \int_0^1 (x - \frac12)^2 = \frac{1}{12}$, $e_2 =
\sqrt{12}\,(X - \tfrac12)$. Then
\[
p_F(f) = \langle f, e_1\rangle e_1 + \langle f, e_2\rangle e_2
= \frac13 + 12\Bigl(\int_0^1 x^2\bigl(x - \tfrac12\bigr)\dd
x\Bigr)\bigl(X - \tfrac12\bigr)
= X - \frac{1}{6},
\]
using $\int_0^1 x^2(x - \frac12)\dd x = \frac14 - \frac16 =
\frac{1}{12}$. The ``least squares'' idea in one line of linear
algebra.
\end{example}

\section{Isometries of the plane}

\begin{definition}\label{def:b1:euclid:isometry}
An endomorphism $u$ of a Euclidean space is an
\emph{isometry}\index{isometry} (or orthogonal map) when it
preserves the norm: $\norm{u(x)} = \norm x$ for all $x$ ---
equivalently (polarization) it preserves the inner product;
equivalently its matrix $A$ in an orthonormal basis satisfies
$A^{\mathsf T} A = I$. Isometries form a group, the orthogonal group
$O(E)$.
\end{definition}

\begin{theorem}[Plane isometries]\label{thm:b1:euclid:plane}
In an orthonormal basis of a Euclidean plane, the matrices of
isometries are exactly
\[
R_\theta = \begin{pmatrix}
\cos\theta & -\sin\theta\\ \sin\theta & \cos\theta
\end{pmatrix}
\quad (\text{rotation of angle } \theta,\ \det = 1),
\]
\[
S_\theta = \begin{pmatrix}
\cos\theta & \sin\theta\\ \sin\theta & -\cos\theta
\end{pmatrix}
\quad (\det = -1),
\]
the latter being the reflection in the line making angle
$\frac\theta2$ with the first basis vector.
\end{theorem}

\begin{proof}
Let $A = \begin{pmatrix} a & c\\ b & d\end{pmatrix}$ with
$A^{\mathsf T}A = I$: columns are unit and orthogonal. The first
column is $(\cos\theta, \sin\theta)$ for some $\theta$; the second,
unit and orthogonal to it, is $\pm(-\sin\theta, \cos\theta)$. The
sign $+$ gives $R_\theta$; the sign $-$ gives $S_\theta$. One checks
$S_\theta^2 = I$ and that the vector $(\cos\frac\theta2,
\sin\frac\theta2)$ is fixed while its orthogonal is reversed: a
reflection. (And $R_\alpha R_\beta = R_{\alpha+\beta}$: the
rotation group is the angle group --- compare
\cref{thm:b1:complex:funceq}.)
\end{proof}

\section{Exercises}

\begin{exercise}[$\star$]\label{exo:b1:euclid:1}
In $\R^3$ canonical: compute $\langle u, v\rangle$, $\norm u$,
$\norm v$ and the angle between $u = (1, 2, 2)$ and $v = (2, -2,
1)$. Verify Cauchy--Schwarz numerically.
\end{exercise}

\begin{exercise}[$\star$]\label{exo:b1:euclid:2}
Prove the parallelogram identity $\norm{x+y}^2 + \norm{x-y}^2 =
2\norm x^2 + 2\norm y^2$ in any Euclidean space, and use it to show
that the sup norm on $\R^2$, $\norm{(x,y)}_\infty = \max(\abs x,
\abs y)$, does not come from an inner product.
\end{exercise}

\begin{exercise}[$\star$]\label{exo:b1:euclid:3}
Apply Gram--Schmidt to $\bigl((1,1,0), (1,0,1), (0,1,1)\bigr)$ in
canonical $\R^3$.
\end{exercise}

\begin{exercise}[$\star$]\label{exo:b1:euclid:4}
In $\R^3$, let $F = \operatorname{Vect}\bigl((1,1,1)\bigr)$.
Determine $F^\perp$ (equation and basis), the matrix of $p_F$ in
the canonical basis, and $d\bigl((1, 2, 3), F\bigr)$.
\end{exercise}

\begin{exercise}[$\star\star$]\label{exo:b1:euclid:5}
(Normal equations) Let $F =
\operatorname{Vect}\bigl((1,0,1),(0,1,1)\bigr) \subseteq \R^3$ and
$x = (1, 1, 4)$. Compute $p_F(x)$ by solving $\langle x - p, v
\rangle = 0$ for the two generators ($p = \alpha(1,0,1) +
\beta(0,1,1)$), then $d(x, F)$. Why is Gram--Schmidt unnecessary
here?
\end{exercise}

\begin{exercise}[$\star\star$]\label{exo:b1:euclid:6}
For $f$ continuous on $\intcc{0}{1}$, prove
\[
\Bigl(\int_0^1 f\Bigr)^{\!2} \leq \int_0^1 f^2 ,
\]
with the case of equality, as an instance of Cauchy--Schwarz in
$C(\intcc{0}{1})$. Then prove $\bigl(\sum_{i=1}^n a_i\bigr)^2 \leq
n \sum a_i^2$ for reals $a_i$.
\end{exercise}

\begin{exercise}[$\star\star$]\label{exo:b1:euclid:7}
Prove that for every subspace $F$ of a Euclidean space:
$(F^{\perp})^{\perp} = F$ and $\dim F^{\perp} = \dim E - \dim F$.
\end{exercise}

\begin{exercise}[$\star\star$]\label{exo:b1:euclid:8}
Identify the plane isometries of matrices
\[
A = \frac{1}{\sqrt 2}\begin{pmatrix} 1 & -1\\ 1 & 1\end{pmatrix},
\qquad
B = \frac{1}{5}\begin{pmatrix} 3 & 4\\ 4 & -3\end{pmatrix}
\]
(type, angle or axis). Compute $A^8$ and $B^2$ without
multiplying matrices.
\end{exercise}

\begin{exercise}[$\star\star\star$]\label{exo:b1:euclid:9}
(Minimum as a projection) Compute
\[
\min_{(a, b) \in \R^2} \int_0^1 \bigl(x^2 - a - bx\bigr)^2 \dd x ,
\]
using \cref{ex:b1:euclid:bestapprox}: the minimum is $\norm{f -
p_F(f)}^2$ for $f = X^2$, $F = \R_1[X]$.
\end{exercise}

\begin{exercise}[$\star\star\star$]\label{exo:b1:euclid:10}
Let $u$ be an isometry of a Euclidean space $E$. Prove that
$\ker(u - \mathrm{id}) \perp \operatorname{im}(u - \mathrm{id})$,
and deduce $E = \ker(u - \mathrm{id}) \oplus \operatorname{im}(u -
\mathrm{id})$. \emph{(Compute $\langle x - u(x), y\rangle$ for
$u(y) = y$, using preservation of the product.)}
\end{exercise}
